\documentclass[11pt]{article}
\usepackage{amsmath,amssymb}
\usepackage[margin=1in]{geometry}
\DeclareMathOperator{\Var}{Var}
\DeclareMathOperator{\pr}{pr}
\DeclareMathOperator{\Cov}{Cov}
\DeclareMathOperator{\AVar}{AVar}

\begin{document}
\centerline{\bf\large Supplementary Material for}
\vspace{2mm}
\centerline{\bf\large ``A comparison index for functional ranking in model selection''}
\vspace{2mm}
\centerline{Yaoping~Wang, Rutgers University}
\vspace{5mm}

\section*{S1.  Nested model derivation for AIC}
\label{sec:aic-detail}

This section provides the derivation underlying the approximation in \S3.5 of the main paper.  Let the full model have $p$ parameters and the restricted model have $s<p$ parameters.  The AIC difference is
\[
T_{f,n} - T_{r,n} = n(\log\hat\sigma_f^2 - \log\hat\sigma_r^2) - 2(p-s),
\]
where $\hat\sigma_k^2$ are the residual variance estimates.

Under the Gaussian likelihood, $n\hat\sigma_r^2$ and $n\hat\sigma_f^2$ are distributed as $\chi^2_{n-p}$ and $\chi^2_{n-s}$ respectively under the null that the restricted model is correct.  The log-likelihood ratio statistic $\ell_n = n\log(\hat\sigma_r^2/\hat\sigma_f^2)$ follows a non-central chi-squared distribution asymptotically:
\[
\ell_n \;\xrightarrow{d}\; \chi^2_{d}(\delta), \qquad d = p-s,
\]
where $\delta$ is the non-centrality parameter.  For fixed local alternatives, $\delta = O(1)$; for fixed alternatives, $\delta$ grows linearly with $n$.

Because the AIC penalty subtracts $2d$, the centred AIC difference is $\ell_n - 2d$.  Writing $\ell_n = \delta + Z$ where $\Var(Z) \approx 2(2\delta + d)$ for the non-central chi-squared distribution, we have
\[
E(T_{f,n} - T_{r,n}) \approx \delta - 2d, \qquad
\Var(T_{f,n} - T_{r,n}) \approx 4(2\delta + d).
\]

The comparison index $\kappa_{fr}$ is therefore
\[
\kappa_{fr} = \frac{|E(T_f - T_r)|}{\{\AVar(T_f - T_r)\}^{1/2}}
           \approx \frac{|\delta - 2d|}{\{4(2\delta + d)\}^{1/2}}
           = \frac{|\delta - 2d|}{2\sqrt{2\delta + d}}.
\]

When $\delta$ is large relative to $d$, the penalty $2d$ is negligible and $\kappa \approx \sqrt{\delta/2}$, giving rapid exponential decay of misranking.  When $\delta \approx 2d$, the numerator vanishes and the functionals become locally indistinguishable: this is the competition regime for AIC-based selection.  Near this point, the misranking rate is $\exp(-n\kappa^2/2) \approx 1$ for moderate $n$, so the selection is essentially random between the two models.

\section*{S2.  Pseudocode for $\hat\kappa$ estimation}
\label{sec:alg}

The plug-in estimator $\hat\kappa$ targets $\sqrt{n_{\text{cal}}}\cdot\kappa$ and is computed from calibration residuals without bootstrap.

\begin{center}
\fbox{\parbox{0.85\textwidth}{
\vspace{2mm}
\textbf{Algorithm:} $\hat\kappa$ estimation from calibration data \\
\vspace{2mm}
\textbf{Input:} calibration set $\{(x_i,y_i)\}_{i=1}^{n_{\text{cal}}}$, fitted estimators $\{\hat f_k\}_{k=1}^m$, level $\alpha$ \\
\textbf{Output:} comparison index $\hat\kappa$ \\
\vspace{2mm}
\begin{enumerate}
  \item[\textbf{1.}] For each estimator $k$, compute calibration residuals $r_{k,i} = |y_i - \hat f_k(x_i)|$.
  \item[\textbf{2.}] Compute CP width: $\hat R_k = 2\cdot Q_{1-\alpha}(\{r_{k,i}\})$, where $Q_{1-\alpha}$ is the empirical quantile.
  \item[\textbf{3.}] Estimate residual scale: $\hat\sigma_k = \operatorname{Std}(\{y_i - \hat f_k(x_i)\})$.
  \item[\textbf{4.}] Normal plug-in density: $\hat f_k = 2\phi(z_{1-\alpha/2})/\hat\sigma_k$, where $\phi$ is the standard normal density.
  \item[\textbf{5.}] Variance estimate: $\widehat{\Var}(\hat R_k) = 4\alpha(1-\alpha)/(n_{\text{cal}}\hat f_k^2)$.
  \item[\textbf{6.}] For each pair $(i,j)$:
    \begin{itemize}
      \item Compute Spearman correlation $\hat\rho_{ij}$ of absolute residual vectors $(r_{i,\cdot}, r_{j,\cdot})$.
      \item $\hat\sigma_{ij} = \{\widehat{\Var}(\hat R_i) + \widehat{\Var}(\hat R_j) - 2\hat\rho_{ij}\sqrt{\widehat{\Var}(\hat R_i)\widehat{\Var}(\hat R_j)}\}^{1/2}$.
    \end{itemize}
  \item[\textbf{7.}] $\hat\kappa = \min_{i<j} |\hat R_i - \hat R_j| / \hat\sigma_{ij}$.
\end{enumerate}
\vspace{2mm}
\textbf{Complexity:} $O(mn_{\text{cal}} + m^2)$.  No bootstrap or cross-validation required.
}}
\end{center}

The variance formula in Step~5 follows from the Bahadur (1966) influence function for the quantile.  The Spearman correlation in Step~6 provides a robust rank-based estimate of the residual dependence between estimators, which suffices because only the sign and approximate magnitude of the correlation affect $\hat\sigma_{ij}$.  The estimator is consistent under the conditions of Assumptions~1--4.

\section*{References}

Bahadur, R. R. (1966). A note on quantiles in large samples. \textit{Ann. Math. Statist.}, 37:577--580.

\end{document}
